% ==============================================================================
% ПАРТИЯ B03-045: Мысли в слух: Часть 2 (блоки DISS-005-B2266 .. DISS-005-B2302)
% Точный перенос рукописи v0.5 без изменений
% ==============================================================================

% DISS-005-TAB-089 (Unwrap figure container с 6 фотографиями)
\begin{center}
% DISS-005-B2266 / DISS-005-FIG-095
\includegraphics[width=0.46\linewidth]{image98.jpeg}\hfill
% DISS-005-B2267 / DISS-005-FIG-096
\includegraphics[width=0.46\linewidth]{image99.jpeg}\par\vspace{0.2cm}
% DISS-005-B2268
\makebox[0.46\linewidth][c]{Фото 1}\hfill
% DISS-005-B2269
\makebox[0.46\linewidth][c]{Фото 2}\par\vspace{0.5cm}

% DISS-005-B2270 / DISS-005-FIG-097
\includegraphics[width=0.46\linewidth]{image100.jpeg}\hfill
% DISS-005-B2271 / DISS-005-FIG-098
\includegraphics[width=0.46\linewidth]{image101.jpeg}\par\vspace{0.2cm}
% DISS-005-B2272
\makebox[0.46\linewidth][c]{Фото 3}\hfill
% DISS-005-B2273
\makebox[0.46\linewidth][c]{Фото 4}\par\vspace{0.5cm}

% DISS-005-B2274 / DISS-005-FIG-099
\includegraphics[width=0.46\linewidth]{image102.jpeg}\hfill
% DISS-005-B2275 / DISS-005-FIG-100
\includegraphics[width=0.46\linewidth]{image103.jpeg}\par\vspace{0.2cm}
% DISS-005-B2276
\makebox[0.46\linewidth][c]{Фото 5}\hfill
% DISS-005-B2277
\makebox[0.46\linewidth][c]{Фото 6}\par\vspace{0.3cm}
% DISS-005-B2278
% DISS-005-B2279
% DISS-005-B2280
% DISS-005-B2281
\end{center}

% DISS-005-B2282
% DISS-005-B2283
% DISS-005-B2284
\par\vspace{0.2cm}

% DISS-005-B2285
\subsection*{Дикие идеи по ИИ}

% DISS-005-B2286
Тут описываются идеи, которые почти гарантированно не воплотятся, но имеют интересные интерпретации и реальную практическую значимость\par\vspace{0.3cm}

% DISS-005-B2287
\subsubsection*{Бот текстового анализа осциллограмм}

% DISS-005-B2288
На основе предыдущих наработок, можно модели проводящей анализ осциллограмм попробовать прилепить какую-то языковую модель, которая будет пытать писать текстовый анализ, просмотрев осциллограмму.\par\vspace{0.3cm}

% DISS-005-B2289
\subsubsection*{Бот техподдержки}

% DISS-005-B2290
Имеется идея создать бота техподдержки, которому будут скормлены знания электроэнергетики и вся наша документация (без конфиденциальной информации). И данный бот сможет отвечать на базовые вопросы.\par\vspace{0.3cm}

% DISS-005-B2291
\par\vspace{0.2cm}
% DISS-005-B2292
\par\vspace{0.2cm}
% DISS-005-B2293
\par\vspace{0.2cm}

% DISS-005-B2294
\section*{БАВР}

% DISS-005-B2295
\subsection*{Задачи}

% DISS-005-B2296
\subsubsection*{Моделирование сети для генерации синтетических данных и более качественной оценки того, как должен работать РНМ}

% DISS-005-B2297
Требуется создать минимум 2 модели сети, которые потом можно и нужно использовать для генерации синтетических данных поведения РНМ. Речь идёт о следующих двух моделях\par\vspace{0.3cm}

% DISS-005-B2298
Первая – это двухсекционная подстанция 6-10~кВ, которая запитана от трансформатора 110/6(10)~кВ, например мощностью 25-63МВА. Суть данной модели в том, чтобы моделировать аварии в сети 110кВ через трансформатор, так как происходят существенные повороты векторов.\par\vspace{0.3cm}

% DISS-005-B2299
Вторая — это двухсекционная подстанция 6-10~кВ, которая запитана от внешней сети этого же напряжения, внешняя сеть относительно большая, там может стоят трансформатор 35/6кВ. Суть данной модели в том, чтобы воспроизвести аварии в аналогичной сети и множество различных случаев ОЗЗ, которые могут перерастать в двухфазные КЗ разных типов, например фаза А во внешней сети и фаза В во внутренней сети. Примечание: Нагрузка на подстанции может быть слабее.\par\vspace{0.3cm}

% DISS-005-B2300
Характер нагрузки на подстанциях можно моделировать одинаковый, в каждом из запусков, прост будет использоваться комбинация их (вкл/откл выключатель или ещё как-то разделённая). На подстанции находится мощная синхронная двигательная нагрузка (около 1МВА), асинхронная нагрузка трансформаторная нагрузка и отходящие линии с средней трансформаторной / двигательной нагрузкой\par\vspace{0.3cm}

% DISS-005-B2301
Так же можно будет замоделировать наш РНМ, для оценки его поведения.\par\vspace{0.3cm}

% DISS-005-B2302
\par\vspace{0.2cm}
